%! AP Statistics — Unit 1, Lesson 1.4 — Warm-Up ANSWER KEY
\documentclass[10pt]{article}
\usepackage{apstats-article}
\usepackage{apstats-key}

%=============================================================================
\begin{document}

\pageheader{Unit 1, Lesson 1.4}{Warm-Up --- Answer Key}
\begin{tabularx}{\linewidth}{X X X}
Name:\;\ans{Key} & Date:\; & Period:\;
\end{tabularx}

\vspace{0.18in}

\noindent Transportation survey ($n = 25$):

\vspace{0.1in}

\begin{center}
\small
\renewcommand{\arraystretch}{1.8}
\begin{tabularx}{0.72\linewidth}{@{} >{\bfseries}p{0.22\linewidth}
                                        C{0.2\linewidth}
                                        C{0.26\linewidth} @{}}
\rowcolor{sky}
\textbf{Mode} & \textbf{Frequency} & \textbf{Rel.\ Frequency} \\
\midrule
Car  & 10 & 0.40 \\
Bus  & 8  & 0.32 \\
Walk & 4  & 0.16 \\
Bike & 3  & 0.12 \\
\midrule
\rowcolor{sky}\textbf{Total} & 25 & 1.00 \\
\end{tabularx}
\end{center}

\vspace{0.18in}

\begin{enumerate}

  \item Most common mode: \ans{Car, used by 40\% of students.}

  \item \ans{Accept any reasonable response: a bar graph with four bars of
        different heights, a circle divided into four sectors, a pictograph,
        etc. The key idea is that a visual display lets you compare sizes at
        a glance without reading numbers.}

  \begin{teachernote}
  The goal here is to surface students' prior knowledge (middle-school bar
  charts) and activate thinking before formal instruction. Do not evaluate
  the sketch quality --- just look for intent to compare categories visually.
  \end{teachernote}

  \item \textit{[Sketch --- see Teacher Note above.]}

  \begin{teachernote}
  A correct sketch would show bars for Car (tallest), Bus, Walk, and Bike
  (shortest) on a labeled axes, OR a pie chart with four sectors of varying
  size. Either is acceptable. Common errors to note for the lesson: missing
  title, y-axis not starting at 0, bars touching.
  \end{teachernote}

  \item \ans{The sketch makes the size comparison between categories faster
        and more intuitive. The table shows the exact numbers (e.g.,\ 0.40
        for car) more precisely. Both tools are useful depending on the
        audience and purpose.}

\end{enumerate}

\vspace{0.14in}

\begin{tcolorbox}[colback=greenbg, colframe=greenacc, boxrule=0.8pt, arc=2mm,
    left=4mm, right=4mm, top=3mm, bottom=3mm]
\small
\begin{itemize}[leftmargin=1.3em, itemsep=5pt]
  \item Today we formalize two visual displays for categorical data:
        \textbf{bar charts} and \textbf{pie charts}.
  \item Both displays use the same relative frequencies from Lesson 1.3.
  \item We will also learn to spot graphs that \textit{lie}.
\end{itemize}
\end{tcolorbox}

\end{document}
